% `tex.web` §90, the tail of `error`, and where its blank line comes from:
%
%     if interaction>batch_mode then decr(selector);
%     ... help lines, each via print_nl ...
%     print_ln;
%     if interaction>batch_mode then incr(selector);
%     print_ln
%
% One `print_ln` closes whatever line the help (or, with `help0`, the error
% context) left open, and the second shows as a blank line. Every report ends
% that way, including the ones that are not faults: §1298 sends `\show`,
% `\showthe`, `\showbox` and `\showlists` through `error` too.
%
% `\show` and `\showthe` were the two reports in this engine that ended without
% it — their `> …` line *is* the message, so they never went through the same
% path as an error with a `!` line above it, and they lost the blank with it.
% trip.log has 67 of these.
%
% §863's paragraph narration is the same fact in another place: `begin_diagnostic
% … end_diagnostic(true)` around the whole of `@firstpass … @@2: line 2.2-`, so
% two consecutive paragraphs are two blocks and not one.
%
% Blank lines are only worth testing because tools/trace_check.sh compares them
% now. It could not before: TeX opens the log with a `(./file.tex` line this
% engine has no equivalent for, and dropping that left a leading blank with
% nothing to match. Leading blanks are dropped on both sides instead, which
% costs nothing — no case begins with one — and leaves every interior blank
% under comparison.
\catcode`\{=1 \catcode`\}=2
\tracingonline=1
\hsize=100pt \parindent=0pt
\font\rip=cmr10 \rip

% The first thing in the log, so that what follows starts at column 0 in both
% engines: TeX's `\message` puts a space in front of itself when the line
% already has something on it, and at the very top of the log that something is
% the `(./file.tex` this engine does not write.
\showthe\hsize
\message{BETWEEN}
\show\hsize
\message{AFTER}

% ...and an error with help text, which ends with the same single blank line
% by a different route: the help's own `print_nl` closes the context line.
\count1=\undefinedzz
\message{END}

% Two paragraphs in a row, so the boundary between their narrations is visible.
% Both break on the first pass on purpose: what the *second* pass narrates is
% its own piece of work and would swamp the blank line this case is about.
\tracingparagraphs=1
AAA AAA AAA AAA\par
AAA AAA AAA AAA\par
\end
